\documentclass[12pt, a4paper]{article}
\usepackage[utf8]{inputenc}
\usepackage{geometry}
\usepackage{listings}
\usepackage{xcolor}
\usepackage{graphicx}
\usepackage{tabularx}
\usepackage{float}
\usepackage{hyperref}

% Page Margins
\geometry{top=1in, bottom=1in, left=1in, right=1in}

% Styling
\definecolor{codegreen}{rgb}{0,0.6,0}
\definecolor{codegray}{rgb}{0.5,0.5,0.5}
\definecolor{codepurple}{rgb}{0.58,0,0.82}
\definecolor{backcolour}{rgb}{0.95,0.95,0.92}

\lstdefinestyle{style}{
    backgroundcolor=\color{backcolour},   
    commentstyle=\color{codegreen},
    keywordstyle=\color{magenta},
    numberstyle=\tiny\color{codegray},
    stringstyle=\color{codepurple},
    basicstyle=\ttfamily\footnotesize,
    breakatwhitespace=false,         
    breaklines=true,                 
    captionpos=b,                    
    keepspaces=true,                 
    numbers=left,                    
    numbersep=5pt,                  
    showspaces=false,                
    showstringspaces=false,
    showtabs=false,                  
    tabsize=2
}

\lstset{style=style}

\title{\textbf{CS327: Compilers} \\ Assignment 3: Syntax Analyzer}
\author{
    \textbf{Group Members:} \\
    Gaurav Budhwani (22110085) \\
    John Twipraham Debbarma (22110108) \\
    Harshit (22110095)
}
\date{\today}

\begin{document}

\maketitle
\tableofcontents
\newpage

\section{Introduction}
This report documents the implementation of the Syntax Analyzer for the \textbf{OPEX} programming language. The parser uses \textbf{Bison} (Yacc) to define and enforce the Context-Free Grammar (CFG), and integrates with \textbf{Flex} for lexical analysis. It successfully parses 12 diverse test cases ranging from basic graph algorithms to computer vision pipelines.

\section{Context-Free Grammar (CFG)}
The following non-exhaustive list of grammar rules illustrates the core structure of the syntax:

\begin{verbatim}
Program     -> PROGRAM ID Declarations MainBlock
MainBlock   -> MAIN BEGIN Statements END | BEGIN Statements END
Declarations-> /* empty */ | Declarations Decl
Decl        -> VarDecl | StructDecl | FuncDecl
VarDecl     -> VAR ID : Type OptSemi 
             | AUTO ID := Expr OptSemi
             | Type ID := Expr OptSemi
Type        -> INT | FLOAT | BOOL | STRING | VOID | TYPE_CV | ArrayType | ID
StructDecl  -> STRUCT ID StructMembers ENDSTRUCT OptSemi
FuncDecl    -> FUNCTION ID ( ParamList ) : Type BEGIN Statements END OptSemi
Statements  -> /* empty */ | Statements Statement
Statement   -> AssignStmt | IfStmt | WhileStmt | ForStmt | ForeachStmt 
             | MatchStmt | IOStmt | ReturnStmt | FuncCallStmt OptSemi
AssignStmt  -> LValue := Expr OptSemi | LValue += Expr OptSemi | LValue ++ OptSemi
IOStmt      -> PRINT Expr OptSemi | INPUT LValue OptSemi
\end{verbatim}

\section{Source Code}
\subsection{parser.y}
The Bison specification that defines the OPEX grammar rules.
\lstinputlisting[language=C, caption=parser.y]{src/parser.y}

\subsection{lexer.l}
The Flex lexer was modified to return Bison tokens instead of printing them out directly.
\lstinputlisting[language=C, caption=lexer.l]{src/lexer.l}

\subsection{Makefile}
The build system compiles both files.
\lstinputlisting[language=make, caption=Makefile]{Makefile}

\section{Test Suite Coverage}
All 12 OPEX test files are available in the \texttt{tests/} directory and pass the parser evaluation correctly:
\begin{itemize}
    \item \texttt{binary\_search.opx}: Array slicing and control flows.
    \item \texttt{edge\_detection.opx}: Custom type instantiation (\texttt{KERNEL}) and implicit declarations (\texttt{AUTO}).
    \item \texttt{sample.opx}: Multi-dimensional arrays and foreach loops.
    \item \texttt{queue.opx} \& \texttt{stack.opx}: Object-struct definitions and usages.
    \item \texttt{floyd\_warshall.opx}: Triply-nested loops and matrix accesses.
    \item ... and six additional algorithmic tests (\texttt{dfs.opx}, \texttt{bst\_inorder.opx}, \texttt{fibonacci.opx}, etc.)
\end{itemize}

\end{document}
